\subsection{The unrestricted Fibonacci-subring case}
\label{sec:unbounded-fibonacci}

\begin{theorem}\label{thm:unbounded-fibonacci}
A rank-four fusion ring containing the Fibonacci fusion ring admits a
pivotal categorification over $\mathbb C$ if and only if it is one of
\[
 K_0(\mathrm{Fib}\boxtimes\operatorname{Vec}(C_2)),\qquad
 K_0(\mathrm{Fib}\boxtimes\mathrm{Fib}),\qquad
 K_0(\operatorname{Ad}(E_8)).
\]
All three admit unitary categorifications. No multiplicity bound is assumed.
\end{theorem}

Write $f^2=1+f$. On the two remaining basis elements, multiplication by
$f$ is a symmetric nonnegative integral matrix $M$ satisfying
$M^2=I+M$. Its eigenvalues are $(1\pm\sqrt5)/2$, so, after reordering,
$M=\left(\begin{smallmatrix}0&1\\1&1\end{smallmatrix}\right)$.
Duality commutes with $M$ and cannot exchange its two basis elements.
Consequently the ring is self-dual and has the following form:
\begin{equation}\label{eq:unbounded-fib-rules}
\begin{gathered}
 f^2=1+f,\qquad fX=Y,\qquad fY=X+Y,\\
 X^2=1+aX+bY,\qquad
 XY=f+bX+(a+b)Y,\\
 Y^2=1+f+(a+b)X+(a+2b)Y,
 \qquad a,b\in\mathbb Z_{\geq0}.
\end{gathered}
\end{equation}
Here rank-four fusion rings are commutative, and the displayed formulas
follow from Frobenius reciprocity and associativity. This two-parameter
family is also identified in \cite[Proposition~3.3]{DongRank4}.
Put
\[
 \phi=\frac{1+\sqrt5}{2},\quad s=a+b\phi,\quad
 x=\frac{s+\sqrt{s^2+4}}2,\quad H=\phi^2,\quad K=\phi+2.
\]
The dimension row is $(1,\phi,x,\phi x)$, and
$D=K(1+x^2)$.

\begin{lemma}\label{lem:unbounded-fib-field}
If $b>0$, the rational fusion algebra in
\eqref{eq:unbounded-fib-rules} is the quartic field
$L=\mathbb Q(\phi,x)$. Pivotal categorifiability is equivalent to
pseudo-unitary categorifiability. A necessary condition for any complex
categorification is
\begin{equation}\label{eq:unbounded-fib-norm}
 N=u^2+uv-v^2=r^2\ \text{or}\ 5r^2,
 \quad u=a^2+b^2+4,\quad v=2ab+b^2,
 \quad r\in\mathbb Z_{>0}.
\end{equation}
\end{lemma}
\begin{proof}
Suppose $x\in\mathbb Q(\phi)$. Both $x$ and $x^{-1}=x-s$ are
algebraic integers, so $\varepsilon=\operatorname{Norm}(x)=\pm1$.
Writing $x'$ for its quadratic conjugate, if $\varepsilon=-1$ then
$s=x+x'$ is rational, contrary to $b>0$. If $\varepsilon=1$, then
$s=x-x'$ has trace zero, contrary to
$\operatorname{Tr}(s)=2a+b>0$. Thus $[L:\mathbb Q]=4$.
Evaluation at the Frobenius--Perron character identifies the
four-dimensional rational fusion algebra with $L$; all its characters
form one Galois orbit. Self-duality makes every pivotal structure
spherical, and \cite[Lemma~\ref{B-lem:galois}]{Finite} gives a pseudo-unitary Galois
conjugate.

Cyclotomicity of Frobenius--Perron dimensions~\cite{ENO} implies that
$L/\mathbb Q$ is abelian, in particular normal. Put
$\Delta=s^2+4=u+v\phi$. The conjugate of $\Delta$ is
$\Delta'=(a+b(1-\phi))^2+4$, so $N=\Delta\Delta'>0$.
Normality of $\mathbb Q(\phi)(\sqrt\Delta)$ implies
$\sqrt{\Delta'}\in L$. Since $\Delta'$ is not a square in
$\mathbb Q(\phi)$, comparison of coefficients in this quadratic
extension shows that $\Delta'/\Delta$ is a square in $\mathbb Q(\phi)$.
Hence $N$ is a square in $\mathbb Q(\phi)$.
A positive rational integer is a square in this quadratic field
precisely when it is an integral square or five times an integral
square, proving \eqref{eq:unbounded-fib-norm}.
\end{proof}

We next determine the twist-one constituents of $I(1)$ without searching
over integral restriction matrices. Write $h=\sqrt N>0$ and
$s'=a+b(1-\phi)$. The four formal codegrees are the four conjugates of
$D$. The involution fixing $\phi$ and sending $x$ to $-x^{-1}$ gives
the ratio $x^2$. For the other two conjugates, direct use of
$x^2=1+sx$ and $\Delta\Delta'=h^2$ gives the ratios
\begin{equation}\label{eq:unbounded-fib-ratios}
 H\left(1+\frac{s-s'\Delta/h}{2}x\right),\qquad
 H\left(1+\frac{s+s'\Delta/h}{2}x\right).
\end{equation}
For example, $K/(3-\phi)=H$ and
$x\sqrt\Delta=1+x^2$ give these identities immediately after
rationalizing $1+y^2$, where $y^2-s'y-1=0$.
The formal-codegree theorem and independence of $1,\phi,x,\phi x$
therefore force the four restriction columns, in that basis order, to be
\begin{equation}\label{eq:unbounded-fib-columns}
 (1,0,0,0),\quad (1,0,a,b),\quad
 (1,1,p,q),\quad(1,1,a+b-p,a+2b-q).
\end{equation}
All entries must be nonnegative integers. In particular, the two
occurrences of $f$ lie one in each of the last two columns; this also
follows from the constant term $H=1+\phi$ in
\eqref{eq:unbounded-fib-ratios}.
For a fully explicit calculation, put
\[
 P=a^3+2a^2b+4a-b^3,\qquad
 Q=a^3+3a^2b+ab^2+4a-2b^3-4b.
\]
These satisfy $Hs'\Delta=P+Q\phi$. If $N=r^2$, define
$(d_p,d_q)=(P/r,Q/r)$; if $N=5r^2$, define
\begin{equation}\label{eq:unbounded-fib-pq}
 (d_p,d_q)=\left(\frac{2Q-P}{5r},\frac{2P+Q}{5r}\right).
 \qquad
 p=\frac{a+b-d_p}{2},\quad q=\frac{a+2b-d_q}{2}.
\end{equation}
The formulas for $p,q$ apply in both cases.

\begin{lemma}\label{lem:unbounded-fib-bound}
If $b>0$ and \eqref{eq:unbounded-fib-rules} admits a pivotal
categorification, then $x<36$, and hence $a+b\phi<36$.
\end{lemma}
\begin{proof}
Pass to a pseudo-unitary categorification by
Lemma~\ref{lem:unbounded-fib-field}. Let $L_X,L_Y$ denote the known
twist-one masses in $I(X),I(Y)$ contributed by the four columns
\eqref{eq:unbounded-fib-columns}. Put
\[
 z_1=p+q\phi,\qquad
 z_2=(a+b-p)+(a+2b-q)\phi,
 \qquad z_1+z_2=Hs.
\]
The last two central dimensions are $H+z_1x,H+z_2x$.
Consequently
\begin{align*}
 L:=L_X+\phi L_Y
 &=sx^2+H(z_1+z_2)+x(z_1^2+z_2^2)\\
 &\geq sx^2+H^2s+\tfrac12H^2s^2x=:L_0.
\end{align*}
\cite[Proposition~\ref{B-prop:optimized-moment}, with the weaker two-moment bound]{Finite}, applied to $X$ and $Y$ and added
with weights $1,\phi$, gives
\[
 9L\leq D\bigl(3(1+\phi^2)x+2(1+\phi)\bigr)
       =K(1+x^2)(3Kx+2H).
\]
Define
\[
 F(t)=\tfrac32(2-\phi)t^2-(8\phi+6)t
                  -\tfrac92(3\phi+6).
\]
Substituting $s=x-x^{-1}$ gives the exact identity
\[
 9L_0-K(1+x^2)(3Kx+2H)
       =\frac{x^2+1}{x}F(x)+\frac{18}{x}.
\]
The left-hand side is nonpositive, so $F(x)<0$.
On the other hand,
\[
 2F(36)=7290-4491\phi>0,\qquad
 F'(36)=210-116\phi>0.
\]
For instance $\phi<34/21$ gives the first strict lower bound $132/7$.
Since $F$ has positive leading coefficient, it is positive for all
$t\geq36$. This proves $x<36$.
\end{proof}

\begin{proof}[Proof of Theorem~\ref{thm:unbounded-fibonacci}]
If $b=0$, the subring on $1,X$ has $X^2=1+aX$.
The rank-two moment obstruction~\cite[Proposition~\ref{B-prop:ranktwo}]{Finite} forces $a=0$ or $1$,
giving the first two rings in the statement. Suppose $b>0$.
Lemmas~\ref{lem:unbounded-fib-field} and
\ref{lem:unbounded-fib-bound} give an exhaustive finite procedure.
The bound $a+b\phi<36$ restricts the parameters to
$0\leq a\leq35$, $1\leq b\leq22$, and is equivalent in that box to
\[
 72-2a-b>0,\qquad 5b^2<(72-2a-b)^2.
\]
For each such pair, compute $u,v,N$ from
\eqref{eq:unbounded-fib-norm} and retain it only if $N$ is a square
or five times a square. Integer square roots test these alternatives
exactly. Next compute the four forced columns in
\eqref{eq:unbounded-fib-columns} and reject a negative or nonintegral
entry. For every remaining pair, compute the masses $L_X,L_Y$ and
test the two-moment inequalities. Each operation takes place in the
algebra defined by $\phi^2=1+\phi$ and $x^2=1+(a+b\phi)x$.

The complete implementation is
\path{companion_software/verify_unbounded_fibonacci.py}.
It exhausts the stated parameter box, constructs the columns and exact
moment deficits, and verifies positive deficits by rational bounds.
The inequalities $3/2<\phi<2$, $x>0$ and $x<a+b\phi+1$ supply its
sign witnesses; the final bound follows from $x=s+x^{-1}$ and $x>1$.
The generated record includes every surviving parameter and its
realization or exclusion certificate. There are no catalogue lookups
or parameter-specific calculations in the argument.
Thus only $(a,b)=(1,1)$ remains when $b>0$. Its fusion rules are those
of $\operatorname{Ad}(E_8)$~\cite{EvansJones}, already identified in the
positive-model table as $F_{3,19}$. The three stated rings have the indicated unitary
models, completing the proof.
\end{proof}
